\documentclass[11pt,a4paper]{article}

\usepackage[utf8]{inputenc}
\usepackage[T1]{fontenc}
\usepackage[spanish,es-noshorthands]{babel}
\usepackage{xcolor}
\usepackage[most]{tcolorbox}
\usepackage{listings}
\usepackage{etoolbox}
\usepackage{fancyhdr}
\usepackage{lastpage}
\usepackage{enumitem}
\usepackage{needspace}
\usepackage[hidelinks]{hyperref}

\setlength{\headheight}{30pt}
\pagestyle{fancy}
\fancyhf{}
\renewcommand{\headrulewidth}{0pt}
\fancyhead[R]{%
  \fontsize{8}{9.6}\selectfont
  Licenciatura en Ciencias de la Computación\\
  Sistemas Operativos\\[2pt]
  \rule{4.2cm}{0.5pt}
}
\fancyfoot[C]{\thepage\ / \pageref{LastPage}}

\newtcolorbox{infobox}{
  colback=blue!6!white, colframe=blue!45!white,
  arc=2mm, boxrule=0.6pt, left=8pt, right=8pt, top=6pt, bottom=6pt,
  title=INFO, fonttitle=\bfseries, coltitle=blue!45!black,
  colbacktitle=blue!6!white, boxsep=1pt, breakable
}

\newtcolorbox{extrabox}{
  colback=orange!8!white, colframe=orange!55!white,
  arc=2mm, boxrule=0.6pt, left=8pt, right=8pt, top=6pt, bottom=6pt,
  title=EXTRA, fonttitle=\bfseries, coltitle=orange!55!black,
  colbacktitle=orange!8!white, boxsep=1pt, breakable
}

\newtcolorbox{notebox}{
  colback=black!5!white, colframe=black!25!white,
  arc=2mm, boxrule=0.5pt, left=8pt, right=8pt, top=6pt, bottom=6pt,
  boxsep=1pt, breakable
}

\lstdefinestyle{codestyle}{
  backgroundcolor=\color{black!5},
  rulecolor=\color{black!25},
  frame=single,
  framerule=0.5pt,
  basicstyle=\ttfamily\small,
  breaklines=true,
  showstringspaces=false,
  tabsize=4,
  keywordstyle=\color{blue!60!black}\bfseries,
  commentstyle=\color{green!40!black},
  stringstyle=\color{orange!70!black},
  columns=fullflexible,
  upquote=true
}
\lstset{style=codestyle}
\BeforeBeginEnvironment{lstlisting}{%
  \par\addvspace{0.2\baselineskip}\noindent
  \begin{minipage}{\linewidth}}
\AfterEndEnvironment{lstlisting}{%
  \end{minipage}\par\addvspace{0.2\baselineskip}}
\newcommand{\inputcode}[2][]{%
  \par\addvspace{0.2\baselineskip}\noindent
  \begin{minipage}{\linewidth}
    \lstinputlisting[#1]{#2}
  \end{minipage}\par\addvspace{0.2\baselineskip}}
\setlist{nosep}

\title{Trabajo Práctico 3 -- Memoria}
\author{}
\date{}

\begin{document}
\maketitle
\thispagestyle{fancy}
\begin{center}
SISTEMAS OPERATIVOS
\end{center}
\newpage

\tableofcontents
\newpage

Este TP sigue el paso de un programa desde el código fuente hasta un proceso
en ejecución. Primero construirás \texttt{hello.c} para observar compilación,
enlace y formatos de archivo. Luego crearás \texttt{memoria.c} para relacionar
sus símbolos, secciones y segmentos con las regiones que el núcleo (kernel)
carga en memoria.

\section{ELF y el proceso de enlace}

ELF (\emph{Executable and Linkable Format}) organiza código, datos, símbolos,
reubicaciones y bibliotecas para que el enlazador y el núcleo puedan usar un
objeto. Para ubicar los tres tipos principales en la ruta del programa,
sigamos el caso de \texttt{hello.c}:

\begin{enumerate}
  \item \textbf{Archivo reubicable (\emph{relocatable file})}. Contiene código
  y datos todavía enlazables con otros objetos. Es la salida de compilar con
  \texttt{gcc -c hello.c}; en Linux suele tener extensión \texttt{.o} (y
  \texttt{.ko} en un módulo del núcleo), y todavía no está listo para iniciar.
  \item \textbf{Archivo ejecutable (\emph{executable file})}. Es el resultado
  usual de enlazar objetos y bibliotecas; queda listo para iniciar. Un script
  de shell no es un ELF: quien se ejecuta es su intérprete. En Linux suele no
  llevar extensión, como el ejecutable \texttt{hello}.
  \item \textbf{Objeto compartido (\emph{shared object})}. Es una biblioteca
  que puede enlazarse o cargarse junto con un programa; en Linux suele tener
  extensión \texttt{.so}. No es un paso obligatorio después del ejecutable,
  sino una rama paralela que puede participar en el enlace o en la carga.
\end{enumerate}

El recorrido principal de este práctico es
\texttt{hello.c} $\longrightarrow$ \texttt{hello.o} $\longrightarrow$
\texttt{hello} $\longrightarrow$ proceso.

Creá \texttt{hello.c} con este contenido:

\inputcode[language=C]{../examples/tp3/hello.c}

\begin{lstlisting}[language=bash]
$ gcc -c hello.c -o hello.o
$ gcc hello.c -o hello
$ file hello hello.o
hello: ELF 64-bit LSB pie executable, x86-64, dynamically linked, ...
hello.o: ELF 64-bit LSB relocatable, x86-64, ...
\end{lstlisting}

\subsection{Del código fuente al ejecutable}

Al compilar, el comando \texttt{gcc} coordina internamente cuatro etapas: el
preprocesador (\emph{preprocessor}) expande \texttt{\#include}, macros y
condicionales; el compilador (\emph{compiler}) traduce el código a ensamblador;
el ensamblador (\emph{assembler}) produce un archivo objeto \texttt{.o}; y el
enlazador (\emph{linker}) combina objetos y bibliotecas para formar el
ejecutable. Una invocación normal ejecuta todas las etapas y entrega el
ejecutable; las opciones siguientes permiten detenerse y conservar resultados
intermedios. Así, \texttt{gcc} coordina el recorrido completo: no hay que
invocar manualmente cada herramienta para compilar un programa habitual:

\begin{lstlisting}[language=bash]
$ gcc -E hello.c -o hello.i
$ gcc -S hello.c -o hello.s
$ gcc -c hello.c -o hello.o
$ gcc hello.o -o hello
\end{lstlisting}

Las extensiones \texttt{.i}, \texttt{.s} y \texttt{.o} identifican,
respectivamente, resultados preprocesados, ensamblador y archivos objeto.
El último comando vuelve a enlazar el objeto y produce el ejecutable.

Una biblioteca estática contiene objetos que se copian al ejecutable; una
compartida queda como dependencia y sus símbolos se resuelven al cargar.
Para comparar ambas variantes usaremos una biblioteca mínima: \texttt{saludo.c}
define la función \texttt{saludar()} y \texttt{main.c} la invoca.

\inputcode[language=C]{../examples/tp3/shared/saludo.c}
\inputcode[language=C]{../examples/tp3/shared/main.c}

Primero compilá los dos archivos fuente como objetos. Este paso es común a las
dos variantes:

\begin{lstlisting}[language=bash]
$ gcc -Wall -Wextra -c saludo.c -o saludo.o
$ gcc -Wall -Wextra -c main.c -o main.o
\end{lstlisting}

\textbf{Enlace estático.} Un archivo \texttt{.a} es normalmente una biblioteca
estática (un archivo \emph{archive}) que agrupa objetos. El prefijo \texttt{lib} y
el sufijo \texttt{.a} permiten que \texttt{-lsaludo} encuentre
\texttt{libsaludo.a}; \texttt{-L.} agrega el directorio actual a la
búsqueda. Creá la biblioteca y enlazala con el programa:

\begin{lstlisting}[language=bash]
$ ar rcs libsaludo.a saludo.o
$ gcc main.o -L. -lsaludo -o saludo-estatico
$ file libsaludo.a saludo-estatico
\end{lstlisting}

\Needspace{12\baselineskip}
\textbf{Enlace dinámico.} Creá un objeto compartido (\texttt{.so}) y
enlazalo dejando la biblioteca como dependencia del ejecutable:

\begin{lstlisting}[language=bash]
$ gcc -Wall -Wextra -fPIC -shared saludo.c -o libsaludo.so
$ gcc -Wall -Wextra main.c -L. -lsaludo -Wl,-rpath,'$ORIGIN' -o saludo-compartido
$ file libsaludo.so saludo-compartido
$ ldd saludo-compartido
$ readelf -d saludo-compartido
\end{lstlisting}

\begin{notebox}
En la variante estática, la implementación de \texttt{saludar} queda dentro
del ejecutable y no hace falta conservar \texttt{libsaludo.a} para
ejecutarlo. En la variante compartida, el ejecutable conserva una
dependencia de \texttt{libsaludo.so}; por eso \texttt{ldd} y
\texttt{readelf -d} permiten verla. Si cambia una biblioteca estática hay
que volver a enlazar; una biblioteca compartida puede reemplazarse sin
recompilar el programa, siempre que se conserve una interfaz compatible.
\end{notebox}

\begin{extrabox}
Compará las variantes con \texttt{file}, \texttt{ldd} y \texttt{readelf -d}:
registrá tamaño, dependencias e intérprete. Estos resultados pueden cambiar
entre arquitecturas y distribuciones.
\end{extrabox}

\subsection{El programa que vamos a observar: \texttt{memoria.c}}

Creá \texttt{memoria.c}. Este único programa contiene los objetos que vamos a
seguir: datos inicializados y no inicializados, una cadena de sólo lectura,
una función, una variable local y una reserva en el heap. También imprime sus
direcciones y su identificador de proceso (PID, \emph{Process ID}), y permanece
vivo para poder observar su mapa.

\inputcode[language=C,basicstyle=\ttfamily\footnotesize]{../examples/tp3/memoria.c}

\begin{lstlisting}[language=bash]
$ gcc -Wall -Wextra -O0 -g -o memoria memoria.c
$ ./memoria
\end{lstlisting}

No copies las direcciones de una ejecución: cambian con ASLR (Address Space
Layout Randomization). En las próximas secciones vamos a comprobar, usando
las herramientas del sistema, qué parte del archivo y qué región del proceso
corresponde a cada objeto.

\section{ELF: cabecera, secciones y símbolos}

El comando \texttt{file} ofrece una clasificación rápida. Para inspeccionar
los campos con más detalle usaremos \texttt{readelf}, una herramienta de
\texttt{binutils} que interpreta las estructuras ELF. Ahora inspeccionamos
el archivo que acabamos de construir, sin cambiar todavía al proceso en
ejecución:

\begin{lstlisting}[language=bash]
$ file memoria
$ readelf -h memoria
$ readelf -h memoria | grep -E 'Class|Data|Type|Machine|Entry'
\end{lstlisting}

La cabecera informa el número mágico (\texttt{7f 45 4c 46}) y la clase del
objeto (\texttt{ELF32} o \texttt{ELF64}). También indica el orden de bytes, la
arquitectura, el tipo, el punto de entrada y la ubicación de las tablas. Con un
ejecutable independiente de posición (PIE, \emph{Position-Independent
Executable}), \texttt{Type} suele ser \texttt{DYN}. Sin PIE puede ser
\texttt{EXEC}; ambos siguen siendo ejecutables. La salida exacta depende de
la arquitectura y de las opciones de compilación.

Usá la salida de tu propio \texttt{memoria}: no copies direcciones de una
salida de referencia. Explicá qué informa cada campo y registrá sus valores.

\subsection{De los objetos fuente a las secciones}

Las secciones son la vista del compilador y del enlazador. Listalas y buscá
los objetos del programa:

\begin{lstlisting}[language=bash]
$ readelf -SW memoria
$ readelf -sW memoria | grep -E 'main|marker|message|global_data|printf'
$ nm -C memoria | grep -E 'main|marker|message|global_data'
$ objdump -d -j .text memoria
$ objdump -s -j .data memoria
$ objdump -s -j .rodata memoria
\end{lstlisting}

\begin{itemize}
  \item \texttt{.text} contiene \texttt{main} y \texttt{marker}; \texttt{.rodata}
  contiene \texttt{message} y otras cadenas.
  \item \texttt{.data} contiene \texttt{global\_data}; \texttt{.bss} (BSS,
  \emph{Block Started by Symbol}) contiene \texttt{global\_data\_2}. El
  archivo guarda el tamaño de \texttt{.bss}, no todos sus ceros.
  \item \texttt{.symtab} contiene símbolos de enlace y depuración, y
  \texttt{.dynsym}, símbolos del enlace dinámico. Un ejecutable \emph{stripped}
  puede perder la primera.
  \item \texttt{local\_data} y \texttt{heap\_data} no son secciones ELF:
  aparecen en la pila y el heap del proceso cuando el programa se ejecuta.
\end{itemize}

En \texttt{readelf -sW}, \texttt{Value} es una dirección o valor relativo y
\texttt{Ndx} identifica la sección. \texttt{printf} suele figurar como
\texttt{UND}, porque su implementación vive en una biblioteca compartida. La
relación entre ese símbolo externo, \texttt{.plt} y \texttt{.got.plt} se
retomará más adelante.

\begin{extrabox}
Copiá \texttt{memoria} como \texttt{memoria-stripped}, ejecutá \texttt{strip}
y repetí \texttt{nm -C} y \texttt{readelf -sW}. Explicá qué información se
pierde y por qué el programa todavía puede ejecutarse.
\end{extrabox}

\section{Segmentos y carga en memoria}

Las secciones organizan el archivo para el compilador y el enlazador; los
segmentos describen lo que el cargador proyecta en memoria. Inspeccioná ambos
vínculos:

\begin{lstlisting}[language=bash]
$ readelf -lW memoria
$ readelf -SW memoria
\end{lstlisting}

La tabla de cabecera de programa (PHT, \emph{Program Header Table}) relaciona
las secciones con los segmentos que el cargador puede proyectar. Cada
segmento \texttt{LOAD} describe una región y sus banderas principales son
\texttt{R} (lectura), \texttt{W} (escritura) y \texttt{E} (ejecución).
Normalmente \texttt{.text} y \texttt{.rodata} van en un segmento \texttt{R E};
\texttt{.data} y \texttt{.bss}, en uno \texttt{R W}.

\texttt{FileSiz} indica cuántos bytes del segmento ocupan espacio en el
archivo; \texttt{MemSiz} indica cuántos bytes necesita en memoria. Para
\texttt{.bss}, el cargador completa con ceros la diferencia entre ambos, por
eso el segmento de datos puede tener \texttt{MemSiz} mayor que
\texttt{FileSiz}. Un segmento puede contener varias secciones. El núcleo
redondea sus límites a páginas, habitualmente de 4 KiB, y conserva sus
protecciones.

\subsection{Seguir un símbolo externo}

Una llamada a \texttt{printf} suele pasar por la PLT (\emph{Procedure Linkage
Table}) y la GOT (\emph{Global Offset Table}): el código salta a la PLT,
ésta usa una entrada de la GOT, el enlazador dinámico resuelve el símbolo y
guarda la dirección para las llamadas siguientes. En la primera llamada, la
entrada de la GOT puede apuntar nuevamente a la PLT; el enlazador dinámico
resuelve la dirección y actualiza esa entrada. En las siguientes llamadas se
usa la dirección ya resuelta.

\begin{lstlisting}[language=bash]
$ objdump -d -j .plt memoria
$ readelf -rW memoria
$ readelf -x .got.plt memoria
$ LD_BIND_NOW=1 ./memoria
$ LD_DEBUG=bindings ./memoria 2>&1 | less
\end{lstlisting}

Las instrucciones concretas dependen de la arquitectura. El enlace perezoso
es una estrategia, no una garantía idéntica entre versiones de PIE y binutils.
Esta vista queda como orientación; el camino principal continúa con la carga
del mismo archivo.

\subsection{Páginas, memoria virtual e intercambio}

La memoria virtual da a cada proceso un espacio propio. El núcleo traduce
direcciones a páginas físicas y las carga bajo demanda: \emph{page-in} trae
una página desde un archivo o intercambio (\emph{swap}); \emph{page-out} la
escribe para liberar memoria. \texttt{vmstat} resume memoria, intercambio,
E/S, sistema y CPU:

\begin{lstlisting}[language=bash]
$ vmstat 1 5
$ free -h
$ cat /proc/meminfo
$ cat /proc/swaps
\end{lstlisting}

\texttt{free} muestra memoria libre; \texttt{si}, páginas llevadas a RAM; y
\texttt{so}, páginas enviadas a swap por segundo. Un \texttt{so} aislado no
diagnostica un problema: correlacioná con \texttt{free}, \texttt{top} y
\texttt{/proc/meminfo}.

\begin{extrabox}
Investigá memoria anónima, páginas respaldadas por archivos y swap en
\texttt{man 5 proc}, especialmente \texttt{meminfo}, \texttt{status} y
\texttt{smaps}. No uses \texttt{swapoff} en una máquina compartida.
\end{extrabox}

\section{\texttt{/proc/<pid>/maps}: del segmento a la región}

El sistema de archivos virtual procfs (\emph{proc filesystem}) expone el
estado del sistema y de los procesos. La información de un proceso con PID
\texttt{1234} está en \texttt{/proc/1234/}; \texttt{/proc/self/} refiere al
proceso que consulta ese archivo.

Ahora ejecutá el mismo \texttt{memoria} y compará sus regiones con los
segmentos de \texttt{readelf -lW}:

\begin{lstlisting}[language=bash]
$ ./memoria &
$ pid=$!
$ cat /proc/$pid/maps
$ cat /proc/$pid/status | grep -E 'Vm|Threads'
$ wait $pid
\end{lstlisting}

Una línea de \texttt{maps} tiene esta forma:

\begin{lstlisting}[language=bash]
55a1...-55a2... r-xp 00002000 08:01 12345 /home/usuario/memoria
7f10...-7f12... r-xp 00000000 08:01 67890 /usr/lib/libc.so.6
7ffd...-7fff... rw-p 00000000 00:00 0      [stack]
\end{lstlisting}

Los campos son: rango virtual (inicio inclusivo, fin exclusivo), permisos
\texttt{rwx}, tipo de compartimiento \texttt{p}/\texttt{s}, desplazamiento
dentro del archivo respaldante, dispositivo, i-nodo y ruta. Los ceros y los
nombres como \texttt{[heap]} o \texttt{[stack]} indican regiones anónimas.
El ejecutable puede aparecer en varias líneas porque sus segmentos tienen
protecciones distintas; también aparecen el enlazador dinámico,
\texttt{libc}, el heap y la pila.


\Needspace{18\baselineskip}
\subsection{Observar el cargador con \texttt{gdb}}

\texttt{gdb} (\emph{GNU Debugger}) es un depurador interactivo: permite
iniciar un programa, detenerlo en un punto elegido e inspeccionar su estado.
Al detenernos en \texttt{main}, el cargador ya preparó el ejecutable y sus
bibliotecas; por eso podemos observar el proceso mientras está detenido.
\texttt{break main} fija el punto de detención, \texttt{run} inicia el
programa, \texttt{backtrace} muestra la pila de llamadas,
\texttt{info proc mappings} muestra sus regiones de memoria,
\texttt{next} avanza una línea y \texttt{continue} reanuda la ejecución.

\begin{lstlisting}[language=bash]
$ gdb ./memoria
(gdb) break main
(gdb) run
(gdb) backtrace
(gdb) print &global_data
(gdb) info proc mappings
(gdb) next
(gdb) continue
(gdb) quit
\end{lstlisting}

También podés usar \texttt{pmap -x <pid>} como vista alternativa. Para
inspeccionar un ELF, \texttt{readelf}, \texttt{objdump} y \texttt{nm}
ofrecen vistas complementarias de sus cabeceras, secciones y símbolos.

\end{document}
